\subsection*{Tile and Flow}

It is possible to use the \texttt{chunkat} expression for either the source or the destination of a \emph{DMA statement}. When \texttt{chunkat} is used only for the source, it typically divides the input or buffer into smaller chunks to move it to storage closer to the processor. In contrast, when \texttt{chunkat} is used only for the destination, it typically moves the buffer to compose part of the output. For example:

\begin{lstlisting}[language=croqtile,caption={Tile and Flow}]
__co__ auto foo(f32 [36] input) {
  f32 [36] output;
  parallel p by 6 : block {
    f0 = dma.copy input.chunkat(p) => shared;  // tile the source
    dma.copy f0.data => output.chunkat(p);     // un-tile the destination
  }
  return output;
}
\end{lstlisting}

In this example, the statement applies \texttt{chunkat} to the destination, moving the data block in \texttt{shared} and "un-tiling" each block to the output. Note that the element types must agree, and that no source dimension may exceed the corresponding destination dimension; a smaller source is accepted.

Using the \texttt{chunkat} expression for both the source and destination is rarely supported by any target platform at this time.
